% Class-compatibility preamble retained for reader excerpts of
%   Hiraku Nakajima, "Quiver Varieties and Branching", SIGMA 5 (2009) 003,
%   DOI 10.3842/SIGMA.2009.003, CC BY 4.0.
%
% The source pins the publisher class sigma, which is not part of the delivery
% baseline, so this file supplies the amsart baseline and then copies the author's
% own declaration block mechanically and unchanged (pinned-original-source.tex
% lines 43-239): the aenume environment, the Theorem/Corollary/Lemma/Proposition and
% definition-style newtheorem declarations, the equation numbering, the author's
% reference macros (\lemref and the like) and the author's mathematical macros.
% Everything from pinned source line 240 onwards is left out: those lines define
% \rf, which uses the publisher class internal dimension \y@boxdim, and the two
% colour rules that serve that one macro; nothing in the delivered excerpts uses
% \rf, \graysquare or \halfgray.
% No mathematical macro, statement, symbol or assumption is introduced, changed or
% removed by this file.  The author's package mathdots (pinned source line 20) is not
% reproduced because it is not installed in the delivery TeX installation and no
% delivered excerpt uses its commands; the author's remaining packages are kept.
\documentclass[10pt,a4paper]{amsart}
\usepackage{amssymb}
\usepackage{amsbsy,amscd,mathtools}
\usepackage{mathrsfs}
\usepackage{nicefrac}
\usepackage[all]{xy}
\usepackage{color}
\usepackage{hyperref}
\newenvironment{aenume}{%
  \begin{enumerate}%
  \renewcommand{\theenumi}{\alph{enumi}}%
  \renewcommand{\labelenumi}{\textup{(\theenumi)}}%
  }{\end{enumerate}}
%

\usepackage{verbatim}
\usepackage{version}
\usepackage{color}
\newenvironment{NB}{
\color{red}{\bf NB}. \footnotesize
}{}
% For a final version, NB disappears.
\excludeversion{NB}
% Older version
%\renewenvironment{NB}{%
%  \comment
%  }{\endcomment}
% The following is used inside NB
\newenvironment{NB2}{
\color{blue}{\bf NB}. \footnotesize
}{}
%       Theorem environments
%% \theoremstyle{plain} %% This is the default
\newtheorem{Theorem}[equation]{Theorem}
\newtheorem{Corollary}[equation]{Corollary}
\newtheorem{Lemma}[equation]{Lemma}
\newtheorem{Proposition}[equation]{Proposition}
%\newtheorem{Axiom}{Axiom}

\theoremstyle{definition}
\newtheorem{Definition}[equation]{Definition}
\newtheorem{Example}[equation]{Example}
\newtheorem{CounterExample}[equation]{Counter-Example}
\newtheorem{Notation}[equation]{Notation}
\newtheorem{Convention}[equation]{Convention}
\newtheorem{Conjecture}[equation]{Conjecture}

%\theoremstyle{remark}
\newtheorem{Remark}[equation]{Remark}
\newtheorem{Remarks}[equation]{Remarks}
%\newtheorem*{Claim}{Claim}
%\renewcommand{\theClaim}{} % to make the claim environment unnumbered
\newtheorem{Question}[equation]{Question}
%\renewcommand{\theQuestion}{}  % to make the notation environment
                               % unnumbered
\newtheorem{Exercise}[equation]{Exercise}
%\renewcommand{\theExercise}{}

\numberwithin{equation}{section}
%\numberwithin{Theorem}{section}
%\numberwithin{figure}{section}

\newcommand{\thmref}[1]{Theorem~\ref{#1}}
\newcommand{\secref}[1]{Section~\ref{#1}}
\newcommand{\lemref}[1]{Lemma~\ref{#1}}
\newcommand{\propref}[1]{Proposition~\ref{#1}}
\newcommand{\corref}[1]{Corollary~\ref{#1}}
\newcommand{\subsecref}[1]{Section~\ref{#1}}
\newcommand{\chapref}[1]{Chapter~\ref{#1}}
\newcommand{\defref}[1]{Definition~\ref{#1}}
\newcommand{\remref}[1]{Remark~\ref{#1}}
\newcommand{\remsref}[1]{Remarks~\ref{#1}}
%
%       Math definitions
%\newcommand{\lsp}[2]{{}^{#1}{#2}}
\newcommand{\lsp}[2]{{\mskip-.3mu}{}^{#1}\mskip-1mu{#2}}
\newcommand{\defeq}{\overset{\operatorname{\scriptstyle def}}{=}}
\newcommand{\interior}{\rule{6pt}{.4pt}\rule{.4pt}{6pt}} % interior
                                                         % product
\newcommand{\C}{{\mathbb C}}
\newcommand{\Z}{{\mathbb Z}}
\newcommand{\Q}{{\mathbb Q}}
\newcommand{\R}{{\mathbb R}}
% Projective spaces
\newcommand{\proj}{{\mathbb P}}
%\newcommand{\CP}{\operatorname{\C P}}
\newcommand{\CP}{\proj}
\newcommand{\RP}{\operatorname{\R P}}
% Lie Groups
\newcommand{\SL}{\operatorname{\rm SL}}
\newcommand{\SU}{\operatorname{\rm SU}}
\newcommand{\GL}{\operatorname{GL}}
\newcommand{\PGL}{\operatorname{PGL}}
\newcommand{\U}{\operatorname{\rm U}}
\newcommand{\SO}{\operatorname{\rm SO}}
\newcommand{\grpSp}{\operatorname{\rm Sp}}
% Lie algebras
\newcommand{\algsl}{\operatorname{\mathfrak{sl}}} % because \sl="slant"
\newcommand{\su}{\operatorname{\mathfrak{su}}}
\newcommand{\gl}{\operatorname{\mathfrak{gl}}}
\newcommand{\algu}{\operatorname{\mathfrak{u}}}
\newcommand{\so}{\operatorname{\mathfrak{so}}}
\newcommand{\algsp}{\operatorname{\mathfrak{sp}}}
\newcommand{\g}{{\mathfrak g}}
\newcommand{\h}{{\mathfrak h}}
% Useful symbols
\newcommand{\Spec}{\operatorname{Spec}\nolimits}
\newcommand{\Proj}{\operatorname{Proj}\nolimits}
\newcommand{\End}{\operatorname{End}}
\newcommand{\Hom}{\operatorname{Hom}}
\newcommand{\Ext}{\operatorname{Ext}}
\newcommand{\Ker}{\operatorname{Ker}}
\newcommand{\Coker}{\operatorname{Coker}}
\newcommand{\Ima}{\operatorname{Im}}
\newcommand{\coker}{\operatorname{coker}}
\newcommand{\Endskew}{\operatorname{Endskew}}
\newcommand{\im}{\mathop{\text{\rm im}}\nolimits}
%\newcommand{\dim}{\mathop{\text{\rm dim}}\nolimits}
\newcommand{\rank}{\operatorname{rank}}
\newcommand{\ind}{\mathop{\text{\rm ind}}\nolimits}
\newcommand{\tr}{\operatorname{tr}}
\newcommand{\Tr}{\operatorname{Tr}}
\newcommand{\DB}{\overline{\partial}}
\newcommand{\pd}[2]{\frac{\partial#1}{\partial#2}}
\newcommand{\deriv}[1]{\frac{d}{d#1}}
\newcommand{\id}{\operatorname{id}}
\newcommand{\ve}{\varepsilon}
%
% notations used in this paper
%
% \newcommand{\vin}[1]{\operatorname{in}(#1)} % incoming vertex
% \newcommand{\vout}[1]{\operatorname{out}(#1)} % outgoing vertex
\newcommand{\vin}[1]{\operatorname{i}(#1)} % incoming vertex
\newcommand{\vout}[1]{\operatorname{o}(#1)} % outgoing vertex
\newcommand{\bM}{{\mathbf M}} % the vector space
\newcommand{\bN}{{\mathbf N}} % the vector space
\newcommand{\M}{{\mathfrak M}} % the moduli space
\newcommand{\Mreg}{\M^{\operatorname{s}}}
\newcommand{\La}{{\mathfrak L}} % the Lagrangian variety
\newcommand{\Fi}{{\mathfrak F}} % the fixed point variety
\newcommand{\dslash}{/\!\!/} % for algebro-geometric quotient (double slash)
\newcommand{\bv}{{\mathbf v}} % the vector v
\newcommand{\bw}{{\mathbf w}} % the vector w
\newcommand{\bC}{{\mathbf C}} % the Cartan matrix
\newcommand{\bA}{{\mathbf A}} % the adjacency matrix
\newcommand{\bI}{{\mathbf I}} % the identity matrix
\newcommand{\bn}{{\mathbf n}} % the kernel of the extended Cartan matrix
\newcommand{\bu}{{\mathbf u}}
\newcommand{\be}{{\mathbf e}}
\newcommand{\bq}{{\mathbf q}}
\newcommand{\bs}{\mathbf s}
\newcommand{\bc}{\mathbf c}
\newcommand{\codim}{\operatorname{codim}} % codimension
\newcommand{\gr}{\operatorname{gr}} % grade
\newcommand{\topdeg}{{\operatorname{top}}} % top degree
\newcommand{\Uq}{{\mathbf U}_q(\mathfrak g)} % the QUE algebra
\newcommand{\Un}{{\mathbf U}_q}
\newcommand{\Um}{\Tilde{\mathbf U}_q(\mathfrak g)} % the modified
                                                 % enveloping algebra
\newcommand{\Ua}{{\mathbf U}_q(\widehat{\mathfrak g})} % the quantum affine
                                                 % algebra
\newcommand{\Uae}{{\mathbf U}_{\varepsilon}(\widehat{\mathfrak g})} %
                                                 % the quantum affine
                                                 % algebra
\newcommand{\Ul}{{\mathbf U}_q({\mathbf L}{\mathfrak g})} % the quantum
                                                      % loop algebra
\newcommand{\Uma}{\Tilde{\mathbf U}^\bw_q(\widehat{\mathfrak g})}
%                                                 the quantum affine algebra
\newcommand{\Ui}{{\mathbf U}^{\Z}_q(\mathfrak g)}
\newcommand{\Ue}{{\mathbf U}_{\varepsilon}(\mathfrak g)}
\newcommand{\Ueone}{{\mathbf U}_{\varepsilon = 1}(\mathfrak g)}
\newcommand{\Uli}{{\mathbf U}^{\Z}_q({\mathbf L}\mathfrak g)}
\newcommand{\Ule}{{\mathbf U}_{\varepsilon}({\mathbf L}\mathfrak g)}
\newcommand{\bU}{\mathbf U} % the enveloping algebra
\newcommand{\Uqm}{{\widetilde{\mathbf U}_q}} % its modified form
\newcommand{\Pa}{{\mathfrak P}} % Hecke correspondence
\newcommand{\HomE}{\operatorname{E}}
\newcommand{\HomL}{\operatorname{L}}
\newcommand{\Zreg}{Z^{\operatorname{reg}\nolimits}}
\newcommand{\hiw}{{L}} % the highest weight module
%      \interval is used to provide better spacing after a [ that
%      is used as a closing delimiter.
\newcommand{\interval}[1]{\mathinner{#1}}
\newcommand{\shfO}{\mathcal O}
\newcommand{\Mp}{\widetilde\M}
\newcommand{\Mpp}{\widehat\M}
\newcommand{\Wedge}{{\textstyle \bigwedge}}
%
\newcommand{\Zw}{Z(\bw)}
\newcommand{\Mw}{\M(\bw)}
\newcommand{\Law}{\La(\bw)}
%\newcommand{\Zw}{Z(\bullet,\bullet;\bw)}
%\newcommand{\Mw}{\M(\bullet,\bw)}
%\newcommand{\Law}{\La(\bullet,\bw)}
\newcommand{\ngpiso}{\overset{.}{=}}
\newcommand{\idl}{\mathfrak I}
\newcommand{\Ktop}{K_{\operatorname{top}}}
\newcommand{\Kztop}{K_{0,\operatorname{top}}}
\newcommand{\Kotop}{K_{1,\operatorname{top}}}

\newcommand{\qH}{\mathbb H}
\newcommand{\D}{\mathcal D}
\newcommand{\A}{\mathcal A}
\newcommand{\incarrow}{\xrightarrow{\mathrm{inclusion}}}
\newcommand{\proarrow}{\xrightarrow{\mathrm{projection}}}
\usepackage{cleveref}
\crefname{originallemma}{original lemma}{original lemmas}
\Crefname{originallemma}{Original lemma}{Original lemmas}
\makeatletter
\newlabel{lem:nondeg}{{2.7}{source}{}{sourcepointer-1}{}}
\newlabel{lem:nondeg@cref}{{[originallemma][0][]2.7}{[1][0][]source}{}{}{}}
\makeatother
\crefname{originallemma}{original lemma}{original lemmas}
\Crefname{originallemma}{Original lemma}{Original lemmas}
\makeatletter
\newlabel{lem:slope}{{2.2}{source}{}{sourcepointer-2}{}}
\newlabel{lem:slope@cref}{{[originallemma][0][]2.2}{[1][0][]source}{}{}{}}
\makeatother
\crefname{originaltheorem}{original theorem}{original theorems}
\Crefname{originaltheorem}{Original theorem}{Original theorems}
\makeatletter
\newlabel{thm:Rudakov}{{2.4}{source}{}{sourcepointer-3}{}}
\newlabel{thm:Rudakov@cref}{{[originaltheorem][0][]2.4}{[1][0][]source}{}{}{}}
\makeatother
\title{Stabilization of the dimension vectors by the stability parameter: the face structure of a zeta-bullet-stable module}
\author{Hiraku Nakajima}
\date{doi:10.3842/SIGMA.2009.003; CC BY 4.0}
\begin{document}
\maketitle
\noindent\textit{Editorial note.} All delivered mathematical text is the author's own, copied from the pinned publisher source sources/algebra-fl01-quiver-varieties-and-branching/pinned-original-source.tex; no statement or proof step was rewritten, reordered or removed.\par\medskip
\noindent\textit{Editorial note.} Preamble provenance: as for the other cards of this source, the retained preamble supplies the amsart baseline in place of the publisher class sigma and copies the author's declaration block of pinned source lines 43-239 unchanged, including the aenume environment and the macros used by the counted proof.\par\medskip
\noindent\textit{Editorial note.} Bibliography provenance: the statement carries the author's own attributions to the references Na-ADHM and Na-quiver, and the proof repeats Na-quiver for the monotonicity statement. They are resolved with sources/algebra-fl01-quiver-varieties-and-branching/bibliography.tex, an excerpt of the author's own thebibliography environment (pinned source lines 4098-4257) copied character for character; no entry was rewritten and no citation key was changed.\par\medskip
\noindent\textit{Editorial note.} External references are navigation pointers only; every deliberately kept step of the author's argument stays inside the counted spans. Their printed numbers are the author's own numbering under the author's newtheorem declarations, computed mechanically by the lane's scanner.\par\medskip
\noindent\textit{Editorial note.} The counted proof is the longest of the paper's second section and contains several displayed computations of the author's own; the delivered span is the author's text from the opening of the proof environment through its closing, with nothing inserted or removed.\par\medskip
\noindent\textit{Allowed standard prerequisites (editorial).}\begin{itemize}
\item The stability parameter zeta-bullet, its faces in the space of stability conditions, and the wall-and-chamber structure attached to the dimension vectors. Developed in the author's subsection on stability (source lines 726-948) and used throughout the counted proof. Supplies the face decomposition with respect to which the module is analysed and the walls along which the dimension vectors stabilize.
\item The monotonicity result for the dimension vectors along a wall, attributed by the author to the cited paper Na-quiver. Used inside the counted proof and attributed by the author to section 5 of Na-quiver. The monotonicity statement is what makes the analysis pass from a general point of the face to the extremal ones.
\item The affine-quiverPlumbing criterion: the author's earlier theorem on which dimension vectors admit a non-empty regular part, used at the end of the argument. The author's own statement at source lines 853-898 (label thm:Rudakov). Identifies the case in which the stabilization claim of the lemma is sharp.
\end{itemize}\medskip
\ifdefined\hypertarget\hypertarget{sourcepointer-1}{}\fi
\noindent\textit{Source reference.} Original lemma 2.7: sources/algebra-fl01-quiver-varieties-and-branching/pinned-original-source.tex, source lines 968--968. Navigation pointer only: the counted proof appeals to this lemma at one place in the induction and the property it supplies is declared as a named standard prerequisite; the lemma's own statement and proof are not reproduced here.\par\medskip
\ifdefined\hypertarget\hypertarget{sourcepointer-2}{}\fi
\noindent\textit{Source reference.} Original lemma 2.2: sources/algebra-fl01-quiver-varieties-and-branching/pinned-original-source.tex, source lines 792--792. Navigation pointer only: it is the author's slope comparison lemma, which the counted proof recalls in one step of the stability analysis; the lemma itself is not reproduced here.\par\medskip
\ifdefined\hypertarget\hypertarget{sourcepointer-3}{}\fi
\noindent\textit{Source reference.} Original theorem 2.4: sources/algebra-fl01-quiver-varieties-and-branching/pinned-original-source.tex, source lines 853--853. Navigation pointer only: the counted proof uses the author's earlier criterion for the non-vanishing of the regular part at the end of the argument, and its content is declared as a named standard prerequisite rather than reproduced.\par\medskip
\begin{Lemma}\label{lem:face}
  Suppose $W\neq 0$.

  \textup{(1) (cf.\ \cite[2.8]{Na-quiver})} Suppose $\zeta$ is in a
  chamber. Then $\zeta$-semistability implies $\zeta$-stability, and hence
  $\Mreg_\zeta(\bv,\bw) = \M_\zeta(\bv,\bw)$.

  \textup{(2) (cf.\ \cite[1.4]{Na-ADHM})} If two stability parameters
  $\zeta$, $\zeta'$ are contained in the same face $F$,
  $\zeta$-stability \textup(resp.\ $\zeta$-semistability\textup) is
  equivalent to $\zeta'$-stability \textup(resp.\
  $\zeta'$-semistability\textup).

  \textup{(3)} Suppose stability parameters $\zeta$, $\zeta^\bullet$ are
  contained in $F$ and $F^\bullet$ respectively. If $F^\bullet\subset
  \overline{F}$, then
  \begin{aenume}\itemsep=0pt
    \item a $\zeta$-semistable module is $\zeta^\bullet$-semistable,
    \item a $\zeta^\bullet$-stable module is $\zeta$-stable.
  \end{aenume}
\end{Lemma}

\begin{proof}
(1) For given $V$,$W$, we def\/ine $\Tilde\zeta=(\zeta,\zeta_\infty)$
with $\theta_{\Tilde\zeta}(V,W) = 0$ as before.
%
Suppose that $(B,a,b)$ is $\Tilde\zeta$-semistable, but not
$\Tilde\zeta$-stable. We take a Jordan--H\"older f\/iltration as in
Theo\-rem~\ref{thm:Rudakov}(2). Since $\delta_k W/\delta_{k+1} W \neq 0$ only
for one $k$ in $0,\dots, N$, we have a $k$ with $\delta_k
W/\delta_{k+1} W = 0$ from the assumption $N\ge 1$. As the quotient
module $\gr_k(B,a,b)\! =\! (B,a,b)|_{(V^k,\delta_k
  W)/\linebreak[2](V^{k+1}, \delta_{k+1} W)}$ is
$\Tilde\zeta$-stable, its dimension vector $\dim V^k - \dim V^{k+1}$
(with the $\infty$-component is $0$) is in $R_+(\bv)$ by
\lemref{lem:nondeg}. Moreover, by the condition in the Jordan--H\"older
f\/iltration, we have
\[
  \theta_{\Tilde\zeta}(\gr_k(B,a,b)) = \theta_{\Tilde\zeta}(V,W).
\]
The right hand side is equal to $0$ by our convention, and the left
hand side is equal to $\zeta\cdot(\dim V^k - \dim V^{k+1})$ up to
scalar. This contradicts with our assumption that $\zeta$ is in a chamber.

(2) We def\/ine $\Tilde\zeta$, $\Tilde\zeta'$ as before.
%
Suppose that $(B,a,b)\in\mu^{-1}(0)\subset \bM(\bv,\bw)$ is
$\Tilde\zeta$-semistable and is not $\Tilde\zeta'$-semistable.
We take the Harder--Narasimhan f\/iltration for $(B,a,b)$ with respect to
the $\Tilde\zeta'$-stability as in \thmref{thm:Rudakov}(1). We have
$N\ge 1$ from the assumption.

Consider f\/irst the case when there exists $0\le k\le k_W - 1$ with
\(
   \theta_{\Tilde\zeta'}(\gr_k(B,a,b)) < 0.
\)
Then
\(
   \theta_{\Tilde\zeta'}(\gr_l(B,a,b)) < 0
\)
for any $0\le l\le k$.
%
On the other hand from the $\zeta$-semistability of $(B,a,b)$, we have
\(
   \theta_{\Tilde\zeta}((V,W)/(V^{k},\delta_{k} W)) \ge 0
\)
by \lemref{lem:slope}.
Therefore there exists at least one $l$ in $[0, k]$ with
\(
   \theta_{\Tilde\zeta}(\gr_l(B,a,b)) \ge 0.
\)
%
We further take a Jordan--H\"older f\/iltration of $\gr_l(B,a,b)$ to f\/ind
$\Tilde\zeta'$-stable representations $V'$ with
\(
  \nicefrac{\zeta'\cdot\dim V'}{\sum \dim V'_i} =
  \theta_{\Tilde\zeta'}(V',0) = \theta_{\Tilde\zeta'}(\gr_l(B,a,b)) < 0.
\)
Note that $\gr_l(B,a,b)$, and hence $V'$ has $0$ in
the $W$-component.
%
Since
\(
   \theta_{\Tilde\zeta}(\gr_l(B,a,b)) \ge 0,
\)
we can take $V'$ so that
\(
    \nicefrac{\zeta\cdot\dim V'}{\sum \dim V'_i}
    = \theta_{\Tilde\zeta}(V',0) \ge 0.
\)
As $\dim V'\in R_+(\bv)$, this contradicts with the assumption that
$\zeta$ and $\zeta'$ are in the common face.

The same argument leads to a contradiction in the case when
there exists $k_W+1 \le k \le N$ with
\(
   \theta_{\Tilde\zeta'}(\gr_k(B,a,b)) > 0.
\)
Since $N\ge 1$ and
\(
   \theta_{\Tilde\zeta'}(V,W) = 0,
\)
at least one of two cases actually occur. Therefore $(B,a,b)$ is
$\zeta'$-semistable.

Suppose further that $(B,a,b)$ is $\zeta$-stable. We want to show that
it is also $\zeta'$-stable. Assume not, and take a Jordan--H\"older
f\/iltration of $(B,a,b)$ with respect to the $\zeta'$-stability. Either
of $\gr_0(B,a,b)$ and $\gr_N(B,a,b)$ has the $W$-component $0$. Suppose
the f\/irst one has the $W$-component $0$. We have
\(
   \nicefrac{\zeta'\cdot \dim (V/V^1)}{\sum \dim V_i/V^1_i} =
   \theta_{\Tilde\zeta'}(\gr_0(B,a,b)) = 0.
\)
On the other hand, the $\zeta$-stability of $(B,a,b)$ implies
\(
   \nicefrac{\zeta\cdot \dim (V/V^1)}{\sum \dim V_i/V^1_i} =
   \theta_{\Tilde\zeta}(\gr_0(B,a,b)) > 0.
\)
This contradicts with the assumption. The same argument applies
to the case when $\gr_N(B,a,b)$ has the $W$-component~$0$. Therefore
$(B,a,b)$ is $\zeta'$-stable.

(3) Take a sequence $\{ \zeta_n \}$ in $F$ converging to $\zeta^\bullet$.
Take a nonzero submodule $(S,\delta W)$ of $(B,a,b)$.
%
(a) From the $\zeta$-semistability and (2), we have
\(
   \theta_{\Tilde\zeta_n}(S,\delta W) \le
   \theta_{\Tilde\zeta_n}(V, W)
\)
for any $n$. Taking limit, we get
\(
   \theta_{\Tilde\zeta^\bullet}(S,\delta W) \le
   \theta_{\Tilde\zeta^\bullet}(V,W).
\)
Therefore $(B,a,b)$ is $\zeta^\bullet$-semistable.
%
(b) We assume $(S,\delta W) \neq (V,W)$.
From the $\zeta^\bullet$-stability, we have
\(
   \theta_{\Tilde\zeta^\bullet}(S,\delta W) <
   \theta_{\Tilde\zeta^\bullet}(V,W).
\)
Then
\(
   \theta_{\Tilde\zeta_n}(S,\delta W) <
   \theta_{\Tilde\zeta_n}(V,W)
\)
for suf\/f\/iciently large $n$. Therefore $(B,a,b)$ is
$\zeta_n$-stable. By (2) it is also $\zeta$-stable.
\end{proof}

% Bibliography excerpt retained for reader excerpts of
%   Hiraku Nakajima, "Quiver Varieties and Branching", SIGMA 5 (2009) 003,
%   DOI 10.3842/SIGMA.2009.003, CC BY 4.0.
%
% Every entry below is copied character for character from the author's own
% thebibliography environment of the pinned publisher source
% (pinned-original-source.tex lines 4098-4257).  No entry was rewritten,
% reordered, merged, split or invented, and no citation key was changed;
% only the entries that the delivered excerpts cite are kept.
\begin{thebibliography}{99}

\bibitem{CB}
Crawley-Boevey W., Geometry of the moment map for representations of
  quivers, {\it Compositio Math.} \textbf{126} (2001),   257--293.
 % \MR{MR1834739 (2002g:16021)}


\bibitem{Kac}
 Kac V.G., Inf\/inite-dimensional Lie algebras, 3rd ed., Cambridge
  University Press, Cambridge, 1990. %\MR{MR1104219 (92k:17038)}


\bibitem{KS}
Kashiwara M., Saito Y., Geometric construction of crystal bases, {\it Duke
  Math. J.} \textbf{89} (1997),  9--36,  \mbox{\href{http://arxiv.org/abs/q-alg/9606009}{q-alg/9606009}}. %\MR{MR1458969 (99e:17025)}


\bibitem{Na-quiver}
Nakajima H., Instantons on {ALE} spaces, quiver varieties, and
  Kac--Moody algebras, {\it Duke Math. J.} \textbf{76} (1994), 365--416.
  %\MR{MR1302318 (95i:53051)}


\bibitem{Na-proc}
Nakajima H., Varieties associated with quivers, in Representation Theory of
  Algebras and Related Topics (Mexico City, 1994), {\it CMS Conf. Proc.}, Vol.~19,
  Amer. Math. Soc., Providence, RI, 1996, 139--157.
  %\MR{MR1388562   (97m:16022)}


\bibitem{Na-alg}
Nakajima H., Quiver varieties and Kac--Moody algebras, {\it Duke Math. J.}
  \textbf{91} (1998), 515--560. %\MR{MR1604167 (99b:17033)}


\bibitem{Lecture}
Nakajima H., Lectures on Hilbert schemes of points on surfaces, {\it University
  Lecture Series}, Vol.~18, American Mathematical Society, Providence, RI, 1999.
  %\MR{MR1711344 (2001b:14007)}


\bibitem{Na-qaff}
Nakajima H., Quiver varieties and f\/inite-dimensional representations of
  quantum af\/f\/ine algebras, {\it J. Amer. Math. Soc.} \textbf{14} (2001),
  145--238, \href{http://arxiv.org/abs/math.QA/9912158}{math.QA/9912158}. %\MR{MR1808477 (2002i:17023)}


\bibitem{Na-tensor}
Nakajima H., Quiver varieties and tensor products, {\it Invent. Math.}
  \textbf{146} (2001), 399--449, \href{http://arxiv.org/abs/math.QA/0103008}{math.QA/0103008}. %\MR{MR1865400 (2003e:17023)}


\bibitem{Na-reflect}
Nakajima H.,  Ref\/lection functors for quiver varieties and {W}eyl group
  actions, {\it Math. Ann.} \textbf{327} (2003), 671--721.
  %\MR{MR2023313   (2004k:16036)}


\bibitem{Na-ADHM}
Nakajima H.,  Sheaves on $ALE$ spaces and quiver varieties, {\it Mosc. Math. J.}
  \textbf{7} (2007),  699--722.


\bibitem{Saito}
Saito Y., Crystal bases and quiver varieties, {\it Math. Ann.} \textbf{324}
  (2002),  675--688, \href{http://arxiv.org/abs/math.QA/0111232}{math.QA/0111232}. %\MR{MR1942245 (2004a:17023)}



\end{thebibliography}
\end{document}
